% Quiz 8: selected shortlist A--D, with simplified numerical answers.
% Solutions: q8-sol-261-II2026.tex. Source notes: q8-shortlist-261-II2026.md.
\documentclass[12pt]{exam}
\usepackage{amsthm}
\usepackage{bm}
\usepackage{libertine}
\usepackage[utf8]{inputenc}
\usepackage[margin=0.5in]{geometry}
\usepackage{amsmath,amssymb}
\usepackage{multicol}
\usepackage[shortlabels]{enumitem}
\usepackage{siunitx}
\usepackage{physics}
\usepackage{booktabs}
\usepackage{graphicx}
\usepackage{pgfplots}
\usepackage{listings}
\usepackage{tikz}
\usepackage{tikz-3dplot}
\usetikzlibrary{arrows.meta,decorations.markings}

\bmdefine{\ii}{i}
\bmdefine{\jj}{j}
\bmdefine{\kk}{k}
\newcommand{\te}{\theta}
\newcommand{\cC}{\mathcal{C}}
\newcommand{\word}[1]{\quad\text{#1}\quad}

\pgfplotsset{width=10cm,compat=1.9}
\usepgfplotslibrary{external}
\tikzexternalize
\tikzexternaldisable

\newcommand{\class}{Math 261-001}
\newcommand{\examnum}{Quiz 8}
\newcommand{\examdate}{October 30}

\begin{document}
\pagestyle{plain}
\thispagestyle{empty}

\noindent
\textbf{\class}\\
\textbf{\examnum}, \textbf{\examdate}

\setlength{\tabcolsep}{3.5cm}
\renewcommand{\arraystretch}{1.5}
\setlength\extrarowheight{1cm}
\begin{tabular}{ |c|c| }
 \hline
 Name   & CSU ID \#  \\
 \hline
\end{tabular}

\vspace{10pt}

Choose and answer \textbf{exactly two} of the four problems. If you attempt
more than two, only the first two will be graded. Write your final answers
in the boxes provided. Only answers in the boxes will be graded.
You may use the surrounding space for scratch work. Each problem is worth
the same number of points.

\begin{enumerate}

\item Find the area of the portion of the surface $z=xy$ satisfying
$x^2+y^2\leq15$.

\vfill
\begin{flushright}
\framebox(250,50){}
\end{flushright}

\item A thin lampshade is formed by rotating the line segment
$y=\frac34x$, $2\leq x\leq4$, about the $y$-axis. Its surface density is
$\delta(x,y,z)=y$. Find the mass of the curved surface, without end caps.

\vfill
\begin{flushright}
\framebox(250,50){}
\end{flushright}

\newpage

\item For each parametrization, write the order of the cross product of the tangent vectors that gives the requested direction for the normal vector. [Hint: Try using the right-hand rule.]

\begin{enumerate}[label=(\alph*),itemsep=4pt]
\item \textbf{Outward} from the ellipsoid:
\[
\vec r(\theta,\phi)
=(2\sin\phi\cos\theta,\ 3\sin\phi\sin\theta,\ \cos\phi),
\qquad0\leq\theta<2\pi,\quad0<\phi<\pi.
\]
\item \textbf{Outward} from the surface of revolution about the $y$-axis:
\[
\vec f(t,\theta)=(\cos\theta,\ t,\ -\sin\theta),
\qquad0\leq\theta<2\pi,\quad0<t<2.
\]
\item \textbf{Upward} on the graph:
\[
\vec r(u,v)=(u,\ v,\ u^2+v^2).
\]
\end{enumerate}

Write each answer as an ordered cross product using the listed parameter
names, such as $\vec r_\theta\times\vec r_\phi$.

\vfill
\begin{flushright}
\begin{minipage}{0.29\linewidth}
\centering
(a)\par\smallskip
\framebox(140,50){}
\end{minipage}\hspace{0.3cm}
\begin{minipage}{0.29\linewidth}
\centering
(b)\par\smallskip
\framebox(140,50){}
\end{minipage}\hspace{0.3cm}
\begin{minipage}{0.29\linewidth}
\centering
(c)\par\smallskip
\framebox(140,50){}
\end{minipage}
\end{flushright}

\item A triangle has vertices
\[
P=(0,0,1),\qquad Q=(2,0,1),\qquad R=(0,3,4).
\]
Consider the parametrizations
\[
\begin{aligned}
\vec a(s,t)&=s\vec{PQ}+t\vec{PR}+P,
&&0\leq s\leq1,\quad0\leq t\leq1,\\
\vec b(s,t)&=t\bigl((s\vec{PQ}+P)-R\bigr)+R,
&&0\leq s\leq1,\quad0\leq t\leq1.
\end{aligned}
\]
\begin{enumerate}[label=(\alph*),itemsep=5pt]
\item Which covers exactly the triangle in question: $\vec a$ or
$\vec b$?
\item For the one that doesn't, which additional inequality on \(s,t\) makes it cover the triangle exactly?
\end{enumerate}

\vfill
\begin{flushright}
\begin{minipage}{0.38\linewidth}
\centering
(a)\par\smallskip
\framebox(180,50){}
\end{minipage}\hspace{0.4cm}
\begin{minipage}{0.46\linewidth}
\centering
(b)\par\smallskip
\framebox(215,50){}
\end{minipage}
\end{flushright}

\end{enumerate}
\end{document}
